\documentclass[10pt,a4paper]{amsart}
\usepackage{amsmath,amssymb,amsthm,graphicx}
% Retained minimal preamble for the Stacks Project chapter more-algebra.tex (More on Algebra).
% The chapter file pinned-original-source.tex carries only \input{preamble} on its first line, so no
% preamble text lives inside the chapter itself.  The Stacks Project uses one single shared preamble for
% every chapter; the snapshot stages it identically as sources/stacks-chow/preamble.tex and
% sources/stacks-etale-cohomology/preamble.tex.  The block below is copied mechanically, line for line,
% from sources/stacks-chow/preamble.tex lines 14-18 (the Xy-pic setup) and 154-217 (the native theorem
% declarations together with the mathematical macros).  Nothing was added, removed or redefined; the
% wrapper class and the amsmath/amssymb/amsthm/graphicx base packages are supplied by the assembler.
% Chapter: more-algebra.tex -- More on Algebra.
% For commutative diagrams we use Xy-pic
\usepackage[all]{xy}

% We use 2cell for 2-commutative diagrams.
\xyoption{2cell}

% Theorem environments.
%
\theoremstyle{plain}
\newtheorem{theorem}[subsection]{Theorem}
\newtheorem{proposition}[subsection]{Proposition}
\newtheorem{lemma}[subsection]{Lemma}

\theoremstyle{definition}
\newtheorem{definition}[subsection]{Definition}
\newtheorem{example}[subsection]{Example}
\newtheorem{exercise}[subsection]{Exercise}
\newtheorem{situation}[subsection]{Situation}

\theoremstyle{remark}
\newtheorem{remark}[subsection]{Remark}
\newtheorem{remarks}[subsection]{Remarks}

\numberwithin{equation}{subsection}

% Macros
%
\def\lim{\mathop{\mathrm{lim}}\nolimits}
\def\colim{\mathop{\mathrm{colim}}\nolimits}
\def\Spec{\mathop{\mathrm{Spec}}}
\def\Hom{\mathop{\mathrm{Hom}}\nolimits}
\def\Ext{\mathop{\mathrm{Ext}}\nolimits}
\def\SheafHom{\mathop{\mathcal{H}\!\mathit{om}}\nolimits}
\def\SheafExt{\mathop{\mathcal{E}\!\mathit{xt}}\nolimits}
\def\Sch{\mathit{Sch}}
\def\Mor{\mathop{\mathrm{Mor}}\nolimits}
\def\Ob{\mathop{\mathrm{Ob}}\nolimits}
\def\Sh{\mathop{\mathit{Sh}}\nolimits}
\def\NL{\mathop{N\!L}\nolimits}
\def\CH{\mathop{\mathrm{CH}}\nolimits}
\def\proetale{{pro\text{-}\acute{e}tale}}
\def\etale{{\acute{e}tale}}
\def\QCoh{\mathit{QCoh}}
\def\Ker{\mathop{\mathrm{Ker}}}
\def\Im{\mathop{\mathrm{Im}}}
\def\Coker{\mathop{\mathrm{Coker}}}
\def\Coim{\mathop{\mathrm{Coim}}}

% Boxtimes
%
\DeclareMathSymbol{\boxtimes}{\mathbin}{AMSa}{"02}

%
% Macros for moduli stacks/spaces
%
\def\QCohstack{\mathcal{QC}\!\mathit{oh}}
\def\Cohstack{\mathcal{C}\!\mathit{oh}}
\def\Spacesstack{\mathcal{S}\!\mathit{paces}}
\def\Quotfunctor{\mathrm{Quot}}
\def\Hilbfunctor{\mathrm{Hilb}}
\def\Curvesstack{\mathcal{C}\!\mathit{urves}}
\def\Polarizedstack{\mathcal{P}\!\mathit{olarized}}
\def\Complexesstack{\mathcal{C}\!\mathit{omplexes}}
% \Pic is the operator that assigns to X its picard group, usage \Pic(X)
% \Picardstack_{X/B} denotes the Picard stack of X over B
% \Picardfunctor_{X/B} denotes the Picard functor of X over B
\def\Pic{\mathop{\mathrm{Pic}}\nolimits}
\def\Picardstack{\mathcal{P}\!\mathit{ic}}
\def\Picardfunctor{\mathrm{Pic}}
\def\Deformationcategory{\mathcal{D}\!\mathit{ef}}
\usepackage{cleveref}
\crefname{originallemma}{original lemma}{original lemmas}
\Crefname{originallemma}{Original lemma}{Original lemmas}
\makeatletter
\newlabel{algebra-lemma-Artin-Rees}{{00IN}{source}{}{sourcepointer-1}{}}
\newlabel{algebra-lemma-Artin-Rees@cref}{{[originallemma][0][]00IN}{[1][0][]source}{}{}{}}
\makeatother
\crefname{originallemma}{original lemma}{original lemmas}
\Crefname{originallemma}{Original lemma}{Original lemmas}
\makeatletter
\newlabel{algebra-lemma-map-AR}{{00IO}{source}{}{sourcepointer-2}{}}
\newlabel{algebra-lemma-map-AR@cref}{{[originallemma][0][]00IO}{[1][0][]source}{}{}{}}
\makeatother
\title{The I-adic homology of a complex of finite modules is pro-isomorphic to the powers of its homology}
\author{The Stacks Project Authors}
\date{stacks-project:more-algebra.tex:local-snapshot-20261012; GNU FDL 1.2 or later}
\begin{document}
\maketitle
\setcounter{subsection}{1}
\noindent\textit{Editorial note.} All delivered mathematical text is the author's own, copied from the pinned Stacks Project chapter file more-algebra.tex; no statement or proof step was rewritten, reordered or removed.\par\medskip
\noindent\textit{Editorial note.} The chapter file pinned-original-source.tex carries only \textbackslash{}input\{preamble\} on its first line, so no preamble text lives inside the chapter. The retained file sources/stacks-more-algebra/retained-preamble.tex is a mechanical, line-for-line copy of the project's single shared preamble as staged in the snapshot; nothing was added, removed or redefined and no package was added.\par\medskip
\noindent\textit{Editorial note.} The source URL is the source-level Stacks Project URL, as required when more than one card is produced from one source; the per-claim tag is carried by claim\_key.\par\medskip
\noindent\textit{Editorial note.} The two cross-chapter references inside the counted proof span (algebra-lemma-map-AR and algebra-lemma-Artin-Rees) are recorded as non-load-bearing navigation pointers; both Artin-Rees applications and everything resting on them stay inside the counted proof span.\par\medskip
\noindent\textit{Allowed standard prerequisites (editorial).}\begin{itemize}
\item The Artin-Rees lemma for an ideal in a Noetherian ring. applied at proof lines 29466-29470 through the cited Algebra lemmas. Supplies the two constants c\_1 and c\_2 whose sum controls the comparison maps.
\item Pro-objects of modules and morphisms between them. the final conclusion at statement lines 29455-29456. Fixes the sense in which (H\_n) and (I\textasciicircum{}nH) are isomorphic.
\end{itemize}\medskip
\ifdefined\hypertarget\hypertarget{sourcepointer-1}{}\fi
\noindent\textit{Source reference.} Original lemma 00IN: focused-seeds/AGP/sources/stacks-snapshot-20261011/tags.txt, source lines 685--685. Cross-chapter navigation pointer into the Stacks Project; the standard fact it supplies is applied inside the counted proof span and the argument resting on it stays there.\par\medskip
\ifdefined\hypertarget\hypertarget{sourcepointer-2}{}\fi
\noindent\textit{Source reference.} Original lemma 00IO: focused-seeds/AGP/sources/stacks-snapshot-20261011/tags.txt, source lines 686--686. Cross-chapter navigation pointer into the Stacks Project; the standard fact it supplies is applied inside the counted proof span and the argument resting on it stays there.\par\medskip
\begin{lemma}
\label{lemma-consequence-Artin-Rees-bis}
Let $I$ be an ideal of a Noetherian ring $A$. Let
$
K \xrightarrow{\alpha} L \xrightarrow{\beta} M
$
be a complex of finite $A$-modules. Set $H = \Ker(\beta)/\Im(\alpha)$.
For $n \geq 0$ let
$$
I^nK \xrightarrow{\alpha_n} I^nL \xrightarrow{\beta_n} I^nM
$$
be the induced complex. Set $H_n = \Ker(\beta_n)/\Im(\alpha_n)$.
Then there are canonical $A$-module maps
$$
\ldots \to H_3 \to H_2 \to H_1 \to H
$$
There exists a $c > 0$ such that for $n \geq c$ the image of $H_n \to H$ is
contained in $I^{n - c}H$ and there is a canonical $A$-module map
$I^nH \to H_{n - c}$ such that the compositions
$$
I^n H \to H_{n - c} \to  I^{n - 2c}H
\quad\text{and}\quad
H_n \to I^{n - c}H \to H_{n - 2c}
$$
are the canonical ones. In particular, the inverse systems
$(H_n)$ and $(I^nH)$ are isomorphic as pro-objects of $\text{Mod}_A$.
\end{lemma}

\begin{proof}
We have $H_n = \Ker(\beta) \cap I^nL/\alpha(I^nK)$.
Since $\Ker(\beta) \cap I^nL \subset \Ker(\beta) \cap I^{n - 1}L$
and $\alpha(I^nK) \subset \alpha(I^{n - 1}K)$ we get the maps
$H_n \to H_{n - 1}$. Similarly for the map $H_1 \to H$.

\medskip\noindent
By Artin-Rees we may choose $c_1, c_2 \geq 0$ such that
$\Im(\alpha) \cap I^nL \subset \alpha(I^{n - c_1}K)$ for $n \geq c_1$ and
$\Ker(\beta) \cap I^nL \subset I^{n - c_2}\Ker(\beta)$ for $n \geq c_2$, see
Algebra, Lemmas \ref{algebra-lemma-map-AR} and
\ref{algebra-lemma-Artin-Rees}. Set $c = c_1 + c_2$.

\medskip\noindent
It follows immediately from our choice of $c \geq c_2$ that for $n \geq c$
the image of $H_n \to H$ is contained in $I^{n - c}H$.

\medskip\noindent
Let $n \geq c$. We define $\psi_n : I^nH \to H_{n - c}$ as follows.
Say $x \in I^nH$. Choose $y \in I^n\Ker(\beta)$ representing $x$.
We set $\psi_n(x)$ equal to the class of $y$ in $H_{n - c}$.
To see this is well defined, suppose we have a second choice $y'$
as above for $x$. Then $y' - y \in \Im(\alpha)$. By our choice of
$c \geq c_1$ we conclude that $y' - y \in \alpha(I^{n - c}K)$
which implies that $y$ and $y'$ represent the same element of $H_{n - c}$.
Thus $\psi_n$ is well defined.

\medskip\noindent
The statements on the compositions $I^n H \to H_{n - c} \to  I^{n - 2c}H$
and $H_n \to I^{n - c}H \to H_{n - 2c}$ follow immediately from our
definitions.
\end{proof}

\end{document}
